% Decoded unit coordinate: a successor to the frozen 111-operation system.
\section{Decoding the unit coordinate}\label{sec:unit-digit}

The addition inserting the constant digit in
$\mathcal C=1+xB+g$ can also be removed. We encode a coordinate
$\delta$ in that position, homogenize the quadratic residuals using
$\delta$, and use a slightly less restrictive mask to force
$\delta=1$. The resulting system has 110 arithmetic instructions.
The encoded coordinates introduced in this construction remain digits
of $g$; they do not increase the number of system unknowns.

\subsection{Homogeneous rows and the unit test}

Take the 1830 integer quadratic residuals $F_i(x,X_1,\ldots,X_{58})$
obtained by the basis selection in Section~\ref{sec:quadratic-masks}.
Put $X_0=x$ and homogenize each residual:
\[
 F_i^{\mathrm h}(\delta,X)
   =\sum_{|I|\le2}a_{i,I}\delta^{2-|I|}X^I.
\]
Introduce one more nonnegative encoded coordinate $u=X_{59}$ and
one ordinary row $u\delta-x^2$. There are now 1831 ordinary homogeneous
quadratic rows. Their target coefficients will be twice their values.
Add a final special target whose coefficient is $\delta^2$.
The new mask will permit a target coefficient to be either zero or one.
Thus each ordinary target, being even, must be zero, while the special
target forces $\delta\in\{0,1\}$. The guard $u\delta-x^2=0$ and the
positive input $x$ then force $\delta=1$.

The distinction in scaling is essential: the special target encodes
$\delta^2$, rather than $2\delta^2$. It is not an additional equation
requiring the special value to vanish.

\subsection{The uniform support and admissible index}

Give $\delta$ weight zero. For the input, the 58 original coordinates,
and $u$, use the 60 positive weights
\[
 v_i=1+3(6963i+i^2)\quad(0\le i\le59),
 \qquad v_*=v_{59}=1242895.
\]
There are 1891 homogeneous quadratic monomials in these coordinates
and $\delta$, and all their weights are distinct. Reduction modulo
three distinguishes the number of positive-weight factors. For two
such factors, division by $6963=2\cdot59^2+1$ recovers the sum and
sum of squares of their indices, hence their unordered pair. This is
the same elementary construction used in the preceding sections.

There are $m=1832$ target positions. Define
\begin{align*}
 D_*&=4v_*+1=4971581,\\
 T_*&=(m+1)D_*+2v_*=9115393763,\\
 t_j&=T_*+jD_*\quad(0\le j<m),\\
 t_*&=18218358574,\qquad K=t_*+1=18218358575.
\end{align*}
The exact checker verifies
\begin{equation}\label{eq:unit-layout}
\begin{gathered}
 D_*>2v_*,\qquad T_*-2v_*>v_*,\qquad
 2T_*-2v_*>t_*,\\
 L=5^{16}=152587890625>3K+2=54655075727.
\end{gathered}
\end{equation}
The 1830 original basis rows, 1832 target positions, and 1891 monomial
weights count different objects.

For each ordinary homogeneous quadratic monomial, let $a$ be its
coefficient in its row and $c\in\{1,2\}$ its multinomial coefficient
in a square. Insert the integer $2a/c$ at the complementary exponent
$t_j-w$, where $w$ is the monomial weight. For the special row insert
the coefficient $1$ at $t_{m-1}$, complementary to the weight zero
of $\delta^2$. Let $\mathcal D(B)$ be the resulting signed coefficient
polynomial. Its row intervals are disjoint by
\eqref{eq:unit-layout}; within a row, weight uniqueness prevents
coefficient collisions.

Choose a power of two $z\ge2$ larger than the absolute values of all
these coefficients, and put $Z=2z$. Define
\begin{align*}
 \ell_0(B)&=1+\sum_{i=1}^{59}B^{v_i}
                  +\sum_{j=0}^{m-1}B^{t_j},\\
 e_0(B)&=z\sum_{h=0}^{K-1}B^h+\mathcal D(B),\\
 V&=\ell_0(Z)+e_0(Z)Z^L.
\end{align*}
The coefficients of $\ell_0$ are zero or one, and those of $e_0$ lie
in $[1,Z-1]$. Both codes have degree $t_*$. The mask includes 1892
non-input coordinates: the unit, 59 other genuine coordinates, and
1832 dummy row coordinates. Including the input gives 1893 coordinates;
we use the convenient strict bound $1900b$ for their coefficient sum.

Choose the fixed power of two $H$ so that
\begin{equation}\label{eq:unit-H}
 H>\max\{2Z^{2L+1},\ 4^{t_*+3}Z\,1900^2,\ 3L,\ 16\}.
\end{equation}
This also gives $H>4Z\,1900^2$ and $H>4Z$.
All index construction depends only on the represented r.e.\ set,
independently of the queried input $x$.

\subsection{The changed system and recovery of the codes}

Retain the reversed-packing system of
Section~\ref{sec:reversed-packing}, including its positive unknown
$\sigma=S_3$. Make only the replacements
\[
 \mathcal C=xB+g,\qquad T_3=(B-4)\ell.
\]
Thus its coding and packing expressions are
\begin{align*}
 B&=Hb^2,& Y&=\ell+eq,\\
 Y+\mathcal C^2+\alpha&=q^2,&Y&=V+t\theta,\\
 S_3&=(2e-Z\lambda)\mathcal C^2+B\lambda(1+q),&\sigma&=S_3,\\
 S&=g+q^2(Y+q^2S_3),\\
 T+1&=q^2(1+\theta\lambda)-(b-1)\ell+(B-4)\ell q^4,\\
 r&=S(n^2-n)+(T+1)(n^2-1).
\end{align*}
All the preceding geometric and Pell equations are retained.

The bootstrap remains valid before any unit coordinate has been
decoded. Indeed $\mathcal C=xB+g>B$, since $x,g>0$. The bound still
gives $\mathcal C<q$ and $q>B$. Replacing $B-2$ by the smaller positive
number $B-4$ preserves the complete bounds $0<S<n$ and $0\le T<n$,
including the treatment of a possibly negative first raw mask.
The preceding borrow argument therefore applies unchanged: it first
recovers the Pell conclusions and then proves that the first-block
borrow is zero. The middle mask decodes
\[
 Y=\ell_0(B)+e_0(B)q,
 \qquad \ell=\ell_0(B)+mq,
 \qquad e=e_0(B)-m,
 \qquad 0\le m<e_0(B).
\]

The low first mask restricts the digits of $g$ to the support of
$\ell_0$, each less than $b$. Its unit digit is some
$0\le\delta<b$. Position one is excluded from $\ell_0$, so the input
digit of $\mathcal C=xB+g$ is exactly $x$. The digit polynomial
$\mathcal C(B)$ has degree at most $t_*$. Write
\[
 \mathcal A=\mathcal C(1)^2,
 \qquad 1\le\mathcal A<1900^2b^2,
 \qquad 4\le Z\mathcal A<B/4.
\]
The lower bound follows from $x>0$ and does not assume $\delta=1$.
Also $e,\mathcal C<B^K$, so $2e\mathcal C^2<2B^{3K}<q$.

To remove the remaining quotient ambiguity, put
\[
 X=(B-4)m=\sum_{j=0}^{t_*+1}X_jB^j,
 \qquad 0\le X_j<B.
\]
The term $(B-4)\ell_0$ is less than $q$. Thus the digits $X_j$
occupy exactly the high third-mask positions $L$ through $L+t_*+1$.
The raw-coefficient and carry calculation in
Section~\ref{sec:reversed-packing} gives digits of $S_3$ at least
$B-Z\mathcal A$ throughout this window. It applies unchanged, because
the formula for $S_3$ is unchanged and its hypotheses were just proved.
The third no-carry condition consequently implies $X_j\le Z\mathcal A-1$.

Consider the digit polynomial $F(U)=\sum_jX_jU^j$.
Since $F(B)=(B-4)m$, the integer $F(4)$ is divisible by $B-4$.
But \eqref{eq:unit-H} gives
\[
 0\le F(4)<Z\mathcal A\,4^{t_*+2}
       <Z\,1900^2b^2\,4^{t_*+2}<B/4<B-4.
\]
Hence $F(4)=0$. Nonnegativity of its coefficients forces every $X_j$
to be zero, so $m=0$. We have recovered both canonical codes.
The evaluation point four is required here by the changed mask;
the earlier evaluation-at-two argument is not used for this step.

\subsection{The relaxed mask forces the unit}

At every tested position, which is below $K$, the residual factor
$e_0-z\lambda$ agrees with $\mathcal D$. The support separation
in \eqref{eq:unit-layout} proves that dummy row coordinates cannot
contribute to any target. Distinct rows cannot interfere, and
quadratic weight uniqueness isolates the intended monomial within
each row. The coefficient of
$\mathcal C^2(e_0-z\lambda)$ is therefore zero at every low mask
position, twice the homogeneous residual at each ordinary target,
and $\delta^2$ at the special target.

Every coefficient of this product has absolute value at most
$z\mathcal A<B/4$. Below position $L$, the offset in $S_3/2$
supplies the digit $B/2$ at every position. The resulting digits
lie strictly between $B/4$ and $3B/4$, without carries between them.
The even third block and mask can both be divided by two. The mask
then has digit
\[
 B/2-2
\]
at each selected position. Its only zero bits within a base-$B$
digit are the unit bit and the highest bit. A digit in
$(B/4,3B/4)$ with no overlap with this mask must consequently be
$B/2$ or $B/2+1$. Its centered coefficient is zero or one.

Every ordinary target coefficient is even, so every ordinary row
vanishes. At the special target, $\delta^2\in\{0,1\}$ gives
$\delta\in\{0,1\}$. The ordinary guard $u\delta-x^2=0$ excludes
$\delta=0$ because $x>0$. Thus $\delta=1$, and the original quadratic
residuals all vanish. This proves soundness without an additional
system equation imposing the unit coordinate.

Conversely, given genuine witnesses, set $\delta=1$ and $u=x^2$,
and set all dummy row coordinates to zero. Choose the power of two
$b$ greater than $x$, $u$, and every original coordinate. The unit
digit guarantees $g>0$. Use the canonical code values. The degree
margin and \eqref{eq:unit-H} give
$Y+\mathcal C^2<q^2$ and $Z\mathcal C^2<q$, hence positive
$\alpha$ and $\sigma=S_3$. The packed quotient is positive, as
before. Ordinary target coefficients are zero and the special
coefficient is one, so all masks hold. The preceding Pell necessity
constructions furnish the remaining positive unknowns.

\begin{theorem}\label{thm:best110}
For each r.e.\ set, the effective admissible index $(Z,V,H)$ above
gives a system in 33 positive unknowns and 21 equations equivalent
to membership of the positive input $x$. Its straight-line certificate
has \textbf{110 arithmetic instructions}: 61 multiplications and
49 additions. Charging the seven fixed-numeral construction instructions
gives 117 operations.
\end{theorem}
\begin{proof}
The equivalence has just been proved. Relative to the 111-instruction
schedule, replacing $1+xB+g$ by $xB+g$ removes one addition.
Replacing $B-2$ by $B-4$ has equal cost, and the numeral $4$ is
already used elsewhere. No positive system unknown or equality is
added. The complete primitive list and all 21 residual identities,
including the existing triangular corrections, are verified by
\file{round8\_1980\_certificate.py}. The fixed exponent remains
$L=5^{16}$, so numeral construction still requires seven instructions.
\end{proof}
